\paragraph{The height at one index.}\label{assembly:wordheight}
Fix an index of the selected family. Its height identity states, for
each word $u$ and for the same constant $C=C_{\rm sz}$, that its
stored height is
\[
 h^{\rm rec}_u=P_u+\sum_{j=1}^m K_{j,u}
             -\sum_{j=1}^m L_{6,j}
             -[K:\mathbb Q]^{-1}\sum_{\sigma:K\hookrightarrow\mathbb C}
                  \log\kappa_\sigma .
\]
Here $P_u=-A\log T_u+e_P^{\max}$ with the record's $C$ in
$e_P^{\max}$. Write $N_j=k+j(t_r+1)$, $N=N_m$, and denote the
record's row weights by $\nu_{j\rho i}(u)$. Then
\begin{align*}
 K_{j,u}
   &=\sum_\rho\sum_{0\le i<N_{j-1}}
       \bigl(\ph(\nu_{j\rho i}(u))-\nu_{j\rho i}(u)\log\Omega\bigr)
       +CfN^2\Lambda_a+Cfm^3N^2L_\varepsilon,\\
 L_{6,j}
   &=\sum_\rho\log H^+_{j\rho}-Cfn_V^2 L_p(b).
\end{align*}
The last expression uses the hook-length product for the positive
partition of Paragraph~\ref{frame:forms}. In factorial coordinates its
logarithm for rows $a_0,\ldots,a_{n-1}$ is
$\sum_{i<n}\log\Gamma(a_i+n-i)-
\sum_{i<l<n}\log(a_i-a_l+l-i)$. The two expressions agree by the
hook-length formula. The mean over the embeddings agrees with the
degree-weighted mean over infinite places when
$\kappa_\sigma=\kappa_{v(\sigma)}$. Substitution proves the height
identity with these inputs at this index.

In $(h^{\rm rec}_u+A\log T_u)-(h^{\rm rec}_{u'}+A\log T_{u'})$
the pairing error, the two uniform kernel errors, the hook terms and
the mean of $\log\kappa$ cancel. In the remaining double sum only
the rows $1\le i<k-1$ can depend on the word. At these rows the two
weights are at least one, differ by at most the small-weight bound
$S$, and are at most $M=6x^6b$ on the selected scale. On $[1,M]$ the
derivative of $\ph$ is $\log y$, bounded by $\log M$. The mean-value
inequality therefore bounds the difference of one summand by
$S(\log M+|\log\Omega|)$. Summing gives
$mfkS(\log M+|\log\Omega|)$. The size estimates on this family give
$kS\le b$ and
$\log M+|\log\Omega|\le(14+6/\varepsilon)L_\varepsilon$.
Consequently the variation is at most
$(20+6/\varepsilon)(N_Ta\Lambda_a+mB_0)$, since
$B_0=fmbL_\varepsilon$ and $m\ge1$. The constant is independent of
the word and scale. This proves (K6) for the selected height and
explains the word index suppressed in $s_j$ when only one word is fixed.

For Definition~\ref{r4:R4:hu} the corresponding term-by-term test is
now explicit. Evaluate its row sums on this realization, padding the
$N_{j-1}$ rows with zeros if needed, use its positive hook factors,
and match its mean of $\log\kappa$. Replacing the sharp errors by
larger errors independent of $u$ adds the same constant to every
$h_u$. It cancels from (K6). In (K4) it adds exactly that constant
to the left-hand side, so its bound must also be included in the
right-hand error. For the frame term this absorption can be made
explicit. Keep the other inputs fixed and replace the per-lift
$Cfn_V^2L_p(b)$ by $C_\lambda fN^2\Lambda_a$ with
$C_\lambda\ge20C$. Paragraph~\ref{frame:forms} makes this a
majorant. Its total increment is at most
$C_\lambda mfN^2\Lambda_a\le 2C_\lambda mfbL_\varepsilon
=2C_\lambda B_0$, using $N^2\le b$ and
Lemma~\ref{r6:R6:lem:sizes}(iii). This is absorbed in the
$m(B_0+N_c)x^{e_D\varepsilon}$ error of (K4) by increasing its
constant by $2C_\lambda$. The height of the independently defined period map
still requires the metric and pairing identifications of
Paragraph~\ref{fr:sec:heights}.
